% 7.1 Compromesso tra remapping L2 continuo e formulazione ALE.
% 7.2 Perché l'ottimizzazione geometrica istantanea non equivale all'ottimo per l'errore PDE temporale.
% 7.3 Generalizzabilità della GNN a diverse condizioni fluidodinamiche senza retraining.
% 7.4 Limitazioni attuali: approccio mono-dominio 2D, connettività costante (r-adaptation pura senza topologia dinamica h/p).

I risultati sperimentali ottenuti attraverso la validazione disaccoppiata (Stage A--C) e il benchmark dinamico end-to-end evidenziano che l'efficacia dell'$r$-adattività neurale in regime transiente dipende strettamente dall'accoppiamento tra generatore geometrico, certificazione di validità elementare e formulazione discreta dell'equazione di evoluzione.

In primo luogo, il confronto tra gli operatori di trasferimento di stato dimostra che il remapping pointwise baricentrico (\textit{GDM + $P_1$}), pur essendo computazionalmente economico ($0.293$ s su $10$ adattamenti), introduce una dissipazione numerica cumulativa che innalza l'errore finale al $12.62\%$ ($12.54\%$ in replay controllato), rendendo l'adattività controproducente rispetto a una griglia fissa ($9.03\%$). L'adozione della proiezione variazionale globale in norma $L^2$ (\textit{GDM + $L^2$}) elimina tale degrado riportando l'errore al $9.11\%$ ($7.97\%$ integrato nel tempo), ma trasferisce l'onere computazionale sulla fase di remapping ($11.153$ s, pari al $90.0\%$ del tempo totale di simulazione).

La formulazione conservativa \textit{GDM + ALE} risolve strutturalmente questo collo di bottiglia: convertendo lo spostamento nodale filtrato $\Delta\mathbf{X}_{\mathrm{safe}}$ in una velocità di griglia continua a tratti e integrando l'operatore cinematico $C_{\mathrm{mesh}}^{n+\frac12}$ conforme alla D-GCL direttamente nel passo temporale, essa preserva esattamente sia il campo costante sia la massa discreta globale ($\mathbf{1}^T M^n\mathbf{U}^n$), abbattendo il tempo totale di esecuzione da $12.387$ s a $2.226$ s (speedup di $5.56\times$ in ciclo chiuso e $5.60\times$ in replay) a fronte di uno scostamento di errore contenuto in $0.34$ punti percentuali rispetto alla proiezione $L^2$.

È opportuno discutere con trasparenza due limitazioni emerse nella presente formulazione, che delineano direzioni naturali di estensione:
\begin{enumerate}
    \item \textbf{Interazione tra errore di discretizzazione temporale e guadagno spaziale}: mentre nello studio statico (Stage A) la rilocazione nodale riduce nettamente l'errore di interpolazione rispetto alla griglia uniforme, nel benchmark dinamico transiente le configurazioni \textit{GDM + $L^2$} ($7.97\%$) e \textit{GDM + ALE} ($8.26\%$) mostrano un errore globale dello stesso ordine della mesh fissa ($7.96\%$), mentre \textit{Oracle + ALE} scende al $7.74\%$. Ciò è dovuto al fatto che, sulla risoluzione $32\times32$ con schema Backward Euler del primo ordine ($\mathcal{O}(\Delta t)$) e assenza di stabilizzazione upwind (es. SUPG), l'errore di dispersione e dissipazione temporale associato all'avanzamento del fronte domina sul guadagno di interpolazione spaziale $\mathcal{O}(h^2)$. L'adozione di integratori temporali ALE di ordine superiore (quali BDF2 o schemi Crank--Nicolson geometricamente conservativi) e l'applicazione a fronti con scale spaziali ancora più ridotte ($\sigma \ll h$) sono attese amplificare significativamente il vantaggio di accuratezza rispetto alla griglia Euleriana fissa.
    \item \textbf{Natura globale del fattore di contrazione $\alpha_{\mathrm{global}}$}: il \textit{Deterministic Safety Filter} garantisce analiticamente $\min_e J_e(\alpha_{\mathrm{global}}) \geq \epsilon$ in tempo $\mathcal{O}(M)$ mediante un unico coefficiente scalare globale. Sebbene nei regimi nominali analizzati la regolarizzazione di training mantenga $\alpha_{\mathrm{global}}=1$, in presenza di griglie di grandi dimensioni un singolo elemento critico prossimo alla degenerazione ($\alpha_e^\star \ll 1$) limiterebbe uniformemente lo spostamento su tutto il dominio. Lo sviluppo di schemi di filtraggio locale con rilassamento ellittico sui soli anelli di adiacenza dell'elemento critico costituisce pertanto un'importante estensione metodologica.
\end{enumerate}
